\section{Related Work}

\subsection{Educational evaluation is broader than answer correctness}

Education-facing evaluation has moved beyond short-answer correctness toward pedagogical quality, adaptivity, and safety. EduBench spans student- and teacher-facing scenarios and scores systems on twelve educational dimensions rather than on answer matching alone \cite{xu2025edubench}. TutorEval focuses on science tutoring over long textbook chapters in open-book and closed-book settings, making it useful for contrasting evidence-grounded tutoring against explanation-first tutoring behaviors \cite{chevalier2024sciencetutors}. MathTutorBench, TutorBench, pedagogical taxonomy work, adaptivity benchmarking, and safety-oriented tutor evaluation reinforce the same methodological point from adjacent angles: a useful tutor must calibrate help, align to learner level, and avoid counterproductive instructional behavior, not just produce a plausible answer \cite{macina2025mathtutorbench,srinivasa2025tutorbench,maurya2025aitutoreval,borchers2025adaptivity,hazra2026safetutors}.

That literature matters for this paper because it changes what counts as a strong architecture. If the target metric is purely factual correctness, retrieval quality may dominate. If the target includes scaffold selection, rubric adherence, or calibrated hinting, orchestration and learner-state reasoning become part of the core problem. Our paper therefore does not introduce another benchmark. Instead, it asks which architectural family best matches the demands that existing educational benchmarks already expose.

\subsection{Classical RAG and agentic RAG}

Classical RAG remains the default architecture for knowledge-grounded assistance because it is simple, inspectable, and often effective on fact-heavy tasks \cite{lewis2020rag}. Subsequent work improves that template through self-reflection, retrieval correction, retrieval-quality diagnostics, and stronger evaluation, as seen in Self-RAG, CRAG, RetrievalQA, RAGChecker, RAGBench, and BERGEN \cite{asai2024selfrag,yan2024crag,zhang2024retrievalqa,ru2024ragchecker,friel2024ragbench,rau2024bergen,chen2023benchmarkingrag}. In those systems, the main problem is still evidence use: how to find better documents, better judge them, and better condition generation on them.

Educational tutoring complicates that picture. Open-book tutoring still requires evidence management, but it also requires deciding when to explain directly, when to ask a guiding question, and when to withhold the final answer to preserve the learning objective. Long-context studies and tutoring benchmarks suggest that access to more context does not automatically solve those interaction problems \cite{liu2024lostmiddle,bai2024longbench2,chevalier2024sciencetutors}. Agentic RAG is therefore relevant here not only as a retrieval recipe, but as evidence that orchestration around retrieval can matter once tasks require planning, staged reasoning, or critique.

\subsection{Multi-agent systems for tutoring and coordination-heavy tasks}

Multi-agent tutoring is attractive because instructional interaction naturally decomposes into diagnosis, planning, tutoring, rubric checking, and critique. Recent systems such as AgentTutor, IntelliCode, and CogEvo-Edu pursue versions of that decomposition, while broader agent frameworks such as ReAct and AutoGen show why explicit roles and actions can outperform a single undifferentiated prompt on harder workflows \cite{liu2026agenttutor,jones2026intellicode,wu2025cogevoedu,yao2023react,wu2024autogen}. At the same time, recent comparative studies caution that coordination overhead can erase gains when the underlying task is narrow or retrieval-dominated \cite{nguyen2025marag,radeva2025coordination}.

The implication for this paper is methodological. ``Agentic'' is too broad a label to support fair comparison. Some systems are essentially classical RAG with better prompting. Others retain retrieval but add explicit planning and critique. Still others use multiple role-specialized agents and invoke retrieval only conditionally. Our taxonomy keeps these families separate so the paper can make narrower claims about when retrieval is the primary bottleneck and when pedagogical coordination is.

\subsection{What current benchmarks can justify}

The paper distinguishes \emph{education-specific} evidence from \emph{transferable} mechanism-level evidence. Education-specific benchmarks such as EduBench, TutorEval, ScienceQA, MathTutorBench, and TutorBench are the only ones that can support claims about pedagogical quality, tutoring usefulness, or adaptive support in educational settings \cite{xu2025edubench,chevalier2024sciencetutors,lu2022scienceqa,macina2025mathtutorbench,srinivasa2025tutorbench}. Transferable benchmarks such as HotpotQA, SCROLLS, LongBench v2, FEVER, Wizard of Wikipedia, BEIR, and AgentBench instead support narrower claims about retrieval depth, evidence justification, grounded dialogue, long-context handling, verification, and coordination \cite{yang2018hotpotqa,shaham2022scrolls,bai2024longbench2,thorne2018fever,dinan2019wizard,thakur2021beir,liu2023agentbench}.

This distinction is central to claim discipline. Strong performance on a general coordination benchmark is useful evidence about orchestration, but it is not sufficient evidence about tutoring quality. Likewise, a tutoring benchmark can support educational claims only up to the limits of its evaluator and task design. The long-form manuscript therefore treats benchmark selection as part of the argument rather than as a background detail.
